\section{Дослідницькі цілі та принципи проєктування}
Аудит операціоналізує шість дослідницьких питань. \textbf{ДП1}: які компоненти джерел та аналізу справді інтегровано і де лишаються ручні інтерфейси. \textbf{ДП2}: які інваріанти походження забезпечуються і чого вони не встановлюють. \textbf{ДП3}: чи відтворюються результати схожості й умовної рандомізації на зафіксованих вхідних даних. \textbf{ДП4} стосується розділення часових позначок комунікацій, наданих інтервалів фізичних подій і невідомого часу. \textbf{ДП5}: які режими семантичних збоїв проходять контрольні бар'єри програмного рівня. \textbf{ДП6}: який обмежений процес роботи аналітика підтримується без припущення про незалежно валідоване розслідування.

Реалізацію організовують три принципи. \emph{Типізоване розділення} тримає окремо джерела, артефакти, твердження, оцінки доказів і висновки. \emph{Явна неповнота} зберігає відсутні вхідні дані, невдалий збір, невідому лінію походження, нерозв'язані події та неперевірені твердження. \emph{Придатні для перегляду обчислення} фіксують вхідні дані, параметри, версії та результати, щоб аналітик міг дослідити, як отримано результат. Ці принципи частково реалізовані, а частково обмежені застарілими інтерфейсами; вони не подаються як універсальні інваріанти для кожного скрипта.

Метод дослідження --- аудит реалізації з обмеженими емпіричними перевірками. Перевагу надано поточному коду та свіжим обчисленням, а не звітам про завершення й планам. Успішний тест встановлює лише перевірену властивість. Кількість вихідних записів встановлює інвентар, а не правильність. Кураторський приклад демонструє робочий процес, а не сліпу точність. Нові публікаційні артефакти зберігають ці розрізнення в матриці «твердження --- докази» та машинозчитуваному маніфесті експериментів. Результати аудиту описують зафіксовану базову ревізію і не змінювалися. Одну оприлюднену післяаудитну зміну правила відношень оцінено окремо в розділі~\ref{sec:rule-v2}, а окремий рукопис кейсу зазнав одного виправлення знаменника.
